\documentclass[11pt]{article}
\usepackage{amsmath,amssymb,amsthm}
\usepackage[margin=1in]{geometry}

\newtheorem{theorem}{Theorem}
\newtheorem{remark}{Remark}

\begin{document}

\title{Formal All-Orders Exceptional-Horizon Frobenius Continuation\\
for the $\ell=2$ Axial GfE System}
\date{6 August 2026}
\maketitle

\paragraph{Scope.}
This note records only the formal local exceptional-horizon theorem
certified symbolically from the canonical GfE axial source. It does not
claim convergence, analytic continuation to the exterior region, a
global black-hole mode, or physical instability.

\begin{theorem}[Formal exceptional-horizon continuation]
Let $M\beta\neq0$ and consider
\[
  \omega=\frac{i}{4M},
  \qquad
  \alpha=\frac12 .
\]
There is a formal logarithmic Frobenius solution of the uneliminated
$\ell=2$ axial GfE equations of the form
\[
  U(x)
  =
  x^{1/2}
  \left[
    u_0+
    \sum_{n\ge1}x^n
      \bigl(u_n+v_n\log x\bigr)
  \right],
  \qquad x=r-2M,
\]
seeded by
\[
  v_1
  =
  c_1
  \begin{pmatrix}
    1\\-2M
  \end{pmatrix},
  \qquad
  u_0
  =
  c_1
  \begin{pmatrix}
    3M/5\\6M^2/5
  \end{pmatrix},
\]
and the total-order-two compatibility relation
\[
  (u_1)_2
  =
  -2M(u_1)_1-\frac{362}{25}Mc_1 .
\]

For every integer $n\ge3$,
\[
A_{nn}
=
\begin{pmatrix}
-\dfrac{\beta^2n(n-1)(2n-3)(2n+1)}{8M^3}
&
 \dfrac{\beta^2n(n-1)(2n-3)}{16M^4}
\\[6pt]
 \dfrac{\beta^2n(n-1)(2n+1)}{16M^4}
&
-\dfrac{\beta^2n(n-1)}{32M^5}
\end{pmatrix}.
\]
It has rank one, with
\[
r_n=
\begin{pmatrix}
1\\2M(2n+1)
\end{pmatrix},
\qquad
\ell_n=
\begin{pmatrix}
1\\2M(2n-3)
\end{pmatrix}.
\]
The adjacent kernel coupling is
\[
\gamma_n
=
\ell_n^{T}A_{n,n-1}r_{n-1}
=
\frac{
2\beta^2(n-1)^2(12n^2-24n-5)
}{
3M^4
}.
\]
Writing $n=m+3$ gives
\[
(n-1)^2(12n^2-24n-5)
=
(m+2)^2(12m^2+48m+31)>0
\]
for every integer $m\ge0$, hence $\gamma_n\neq0$.

The projected logarithmic equation has the form
\[
  \ell_n^T R_n+\gamma_n\kappa_{n-1}=0,
\]
and therefore uniquely fixes $\kappa_{n-1}$ while leaving a new
current kernel coefficient $\kappa_n$.

For the non-logarithmic equation, the exponent-derivative source
$D_{nn}$ satisfies
\[
  \ell_n^T D_{nn}r_n=0.
\]
Thus the free current logarithmic kernel does not enter the
non-logarithmic compatibility condition
\[
  \ell_n^T T_n+\gamma_n\mu_{n-1}=0.
\]
This uniquely fixes $\mu_{n-1}$ while leaving a new current ordinary
kernel coefficient $\mu_n$.

Hence the logarithmic and non-logarithmic recurrences advance
inductively to every Laurent order, with no finite-order algebraic
compatibility obstruction on this exceptional local horizon branch.
\end{theorem}

\begin{remark}[Exact boundary]
The theorem is formal and local. It does not establish coefficient
growth bounds, convergence of the infinite Frobenius series, analytic
continuation away from the horizon, global black-hole boundary
conditions, or physical instability.
\end{remark}

\end{document}
